\documentclass{article}
\usepackage{amsmath,amssymb}
\emergencystretch=4em
\begin{document}
\section*{Regional coordinated traffic signal control: A complete definition of the current problem.}
This problem involves optimizing regional coordinated traffic signal control strategies, using the given intermediate road network, vehicle demands, and signal phases as input parameters.

\subsection*{Environment and status}
The directed road network and agent regions are provided as input. The road segment $e$ has a maximum storage capacity of $C_e$, number of lanes $n_e$, length $L_e$, free driving time $\tau_e$, and green light phase number. The phases are predefined as four mutually exclusive directions. The complete environmental state $X_t$ includes information such as the FIFO queues of vehicles at each road segment, vehicles in motion, remaining driving time, the vehicle’s shortest possible path, path location, queues at starting points outside the network, current phase, start time of the green light, countdown timer for yellow/red lights, and the clock. This complete state includes Markov transition probabilities. The following observations represent a compressed representation of the complete state.

The starting point $o$ generates Poisson arrivals per second, with a rate of
\[
\lambda_o(t)=\lambda_0m_o
\frac{1+0.23\sin(2\pi t/H+0.27o)+0.13\sin(4\pi t/H+0.6)}
{H^{-1}\sum_{u=0}^{H-1}[1+0.23\sin(2\pi u/H+0.27o)+0.13\sin(4\pi u/H+0.6)]}.
\]
Destinations are drawn uniformly from the remaining nodes, and each route follows a shortest path by link length. Vehicles unable to enter the network remain in the origin FIFO queue. Each second, the system generates demand, admits vehicles subject to first-link capacity, places vehicles that finish traversing a link in that link's FIFO queue, updates the signals, and processes intersection departures in rotating order $(t+j)\bmod |V|$. A green link discharges at most one vehicle per lane every two seconds. Vehicles are held when downstream storage is insufficient. A vehicle released from a link either enters the next link and travels for $\tau_e$ seconds or completes its trip. The following conditions must hold throughout:
\[
0\le O_e(t)\le C_e,\quad O_e(t)=Q_e(t)+M_e(t),\qquad
G(t)=D(t)+\sum_eO_e(t)+B(t).
\]
$Q,M,O$ represent the number of vehicles in queue, in transit, and occupied, while $G,D,B$ represents the cumulative number of trips generated, completed, and backlog at the starting point. All unfinished trips and backlog at the starting point should be reported.

\subsection*{Strategies and feasibility of signals}
Every $\delta$ seconds, select action $a_i\in\{0,1\}$ for each area. Action 0 operates according to $\arg\max_{e\in\operatorname{in}(v)}O_e/(L_en_e)$, selecting the preset green light phase. This corresponds to the density and preset phase options defined in Algorithm 1. Action 1 rotates through the phases sequentially over time. The switch only occurs when the current green light duration reaches its minimum value; when the maximum green light duration is reached, it must be switched to another phase. Yellow lights and full red lights remain activated for a specified number of seconds, during which no vehicles are allowed to pass. After this period, the new green light starts counting from zero again. The last segment of the green light at the end of the simulation can be shorter than the minimum duration, but must not exceed the maximum duration.

The observations of the agent $i$ consist of six items: the total number of queues, categorized by level 3 thresholds (5, 10), the IDs of the two most congested intersections, information about the classification of the two intersections, and the result of the previous action taken. These observations are combined with those of neighboring agents, based on their distance from the region’s center. The average action of neighboring agents $\bar a_i=|\mathcal N_i|^{-1}\sum_{j\in\mathcal N_i}a_j$ is used in the implementation; the sum of actions from neighboring agents serves as an equivalent index, since the number of neighbors per region is fixed. The selection of candidate strategies is based on the current observations and historical action signals.

\subsection*{Objectives and solving algorithms}
For the actually evaluated $E$ independent demand sample trajectories, the goal is to…
\[
\min_\pi J_E(\pi)=\frac1{EAT}\sum_{e=1}^E\sum_{t=1}^T
\sum_{i=1}^A\sum_{v\in G_i}\sum_{l\in\operatorname{in}(v)}Q_l^{e,t}.
\]
This represents the actual average number of waiting passengers along the route, based on recorded data. Waiting time at the starting point, which is not part of this metric, will be reported separately. The solution employs Q-learning as a heuristic method to find feasible strategies. The reward mechanism is as follows:
\[
f_i^t=\frac1{|G_i|}\sum_{v\in G_i}
\frac{\sum_{l\in\operatorname{in}(v)}M_l^t v_l}{\sum_{l\in\operatorname{in}(v)}L_ln_l},
\quad r_i^t=\operatorname{sgn}(f_i^t-f_i^{t-1})
+\sum_{j\in\mathcal N_i}\frac{\operatorname{sgn}(f_j^t-f_j^{t-1})}{d_{ij}}.
\]
\[
Q_i\leftarrow(1-\alpha)Q_i+\alpha
\left[r_i+\gamma\max_{a'}Q_i(X',a',\bar a_i')\right],\qquad
\pi_i(a)\propto\exp(+\beta Q_i(a)).
\]
The “termination step” removes any future-related items; the positive sign corresponds to the direction that maximizes rewards. This approach ensures that the logit values for two actions are limited to $[-60,60]$ to avoid floating-point overflow issues. Actual probabilities are determined based on these limited values. Each scenario has a budget of 3600 seconds, and time limits are checked according to the control steps. The number of training iterations and evaluation sessions are determined by the parameters set for the algorithm. The output includes feasible control strategies and their corresponding trajectory objectives.
\end{document}
